%----------------------------------------------------------------------------
\chapter{\bevezetes}
%----------------------------------------------------------------------------

\label{chap:introduction}

%----------------------------------------------------------------------------
\section{Motivation}
\label{sec:motivation}
%----------------------------------------------------------------------------

In many graph datasets, nodes have textual descriptions. For example, the nodes of a citation network represent papers, each associated with an abstract. When a new node is added to such a graph, its relationships to existing nodes may already be known while its own text is not yet available. The existing neighbours, in contrast, already have textual descriptions. Graphs whose nodes carry text in this way are called text-attributed graphs (\autoref{sec:tag}). This situation raises the question of whether the relationships of the new node, together with the descriptions of its neighbours, provide enough context to generate useful text for it.

The graph identifies which existing nodes are connected to the new node, while their textual descriptions provide information about what those neighbours represent. Together they indicate the new node's likely topic, role or characteristics, and a language model could use this context to write a description. This thesis investigates how this neighbourhood information can be encoded into a soft prompt: a sequence of continuous vectors that conditions a language model to generate text for the new node (\autoref{sec:soft-prompts}).

Citation graphs make this question concrete. A newly considered paper has references that connect it to existing publications. These publications indicate the research the paper builds on, and because papers often cite related research~\cite{zhu2020beyond}, they give context for its topic and objectives. They do not, however, uniquely determine its particular contribution, methods or findings: two papers can build on similar references and still contribute different things. Because published papers already have abstracts, citation graphs also provide target texts against which generated abstracts can be evaluated.

%----------------------------------------------------------------------------
\section{Research question and goal}
\label{sec:research-question}
%----------------------------------------------------------------------------

The central question of this thesis is whether a language model can generate the text of a new node from its known graph neighbourhood through soft prompting. In the citation graph studied here, the target text is a paper's abstract, and the neighbourhood consists of the papers it cites, its references. The titles and abstracts of the references are first encoded as embeddings, numerical vectors representing their content, so the language model does not read them as text. The target paper's own title, abstract and embedding are not part of the generation input. Papers that later cite the target are also excluded, because they may contain information that was not available when its abstract was written. \autoref{chap:task} defines the task and these restrictions precisely.

I evaluate this task on an existing citation graph, withholding each target paper's own text from the generation input and comparing the generated abstract with the original. The goal is agreement with the target paper's specific content rather than reproduction of its wording. This requires distinguishing abstracts that recover target-specific content from fluent but generic scientific prose, which can read well without describing the target paper.

%----------------------------------------------------------------------------
\section{Approach and related work}
\label{sec:approach}
%----------------------------------------------------------------------------

Language models have already been combined with graph encoders. G-Retriever~\cite{he2024gretriever}, for example, answers questions about graphs whose nodes and edges carry text: a graph neural network (GNN) encodes a selected subgraph into a vector that the language model receives as a soft prompt, and the same subgraph is also supplied to the model as text. This thesis instead investigates generating a target node's own text from its neighbourhood, with paper-specific information reaching the language model through continuous representations. \autoref{sec:related-work} discusses this approach in more detail.

The implemented model follows the flow of information from the references to the generated abstract. A pre-trained embedding model encodes the title and abstract of each reference into an embedding. A GNN combines the embeddings of a target paper's references and estimates the embedding of the target paper itself: the graph determines which papers contribute, and their embeddings supply information about what those papers describe. A prefix adapter converts this estimate into attention prefixes, continuous vectors that the language model attends to while it generates. The language model, Qwen3-14B~\cite{yang2025qwen3}, then generates the abstract. It is adapted with low-rank adaptation (LoRA)~\cite{hu2022lora}, which adds small trainable matrices while the original weights of the model remain frozen. \autoref{chap:method} describes the model in detail.

The prefix adapter and the adapted language model together form the reader. I first train the reader and GNN estimator separately, then fine-tune them jointly. To build the dataset, I enriched the ogbn-arxiv citation graph~\cite{hu2020open} with citation records from Semantic Scholar~\cite{kinney2023semantic}. I implemented the full pipeline and evaluated the generated abstracts with text-similarity measures, a word-based identification test and content agreement assessed by a language-model judge. Ablations compare the learned estimate with a parameter-free average of the reference embeddings and use two diagnostic controls: an input without paper-specific information and the target paper's own embedding.

%----------------------------------------------------------------------------
\section{Scope}
\label{sec:scope}
%----------------------------------------------------------------------------

The implementation makes several specific design choices within the assigned task. The GNN estimate conditions the language model through key and value prefixes in every attention layer rather than through vectors projected into its input embedding space. The reader and GNN estimator are initially trained separately, followed by joint fine-tuning of all trainable components. The structure-agnostic baseline is the parameter-free reference mean. The reasons for these choices are discussed in \autoref{chap:discussion}. The evaluation concerns one enriched citation graph.

The work includes enriching ogbn-arxiv with Semantic Scholar citation records, training an estimator of the target paper's text embedding, and checking the behaviour of the evaluation measures. The implementation will be published on GitHub: \textcolor{red}{[TODO: repository URL]}.

%----------------------------------------------------------------------------
\section{Main findings}
\label{sec:findings}
%----------------------------------------------------------------------------

The experiments show that neighbourhood-derived soft prompts can guide a language model to generate coherent scientific abstracts. In successful cases, the generated texts capture the target paper's research area, problem and task, and contain vocabulary that distinguishes the target paper from other papers. The learned GNN estimator improves content agreement over averaging the reference embeddings. The evaluation, however, assesses the research context rather than the target paper's methods or findings, and the examples suggest that the neighbourhood guides the former more than the latter. A generated abstract can therefore describe a plausible study in the relevant area, but should not be treated as a faithful reconstruction of how the target study was conducted or what it found.

%----------------------------------------------------------------------------
\section{Outline of the thesis}
\label{sec:outline}
%----------------------------------------------------------------------------

\autoref{chap:background} introduces the technical background and related work, and \autoref{chap:task} formally defines the task and the information available to the model. \autoref{chap:data} and \autoref{chap:method} describe the dataset preparation, the model architecture and the training procedure.

\autoref{chap:evaluation} introduces the evaluation methodology. \autoref{chap:results} reports the training and the evaluation of the final model, and \autoref{chap:ablations} compares the configurations to show what each component contributes. \autoref{chap:discussion} interprets the results and explains the design choices, and \autoref{chap:conclusion} summarises the work and outlines future work. The appendices contain supporting material, including the judge prompts and example generations.
